\documentclass[11pt]{article}
\usepackage[margin=1in]{geometry}
\usepackage{amsmath,amssymb,amsthm,mathtools}
\usepackage[hidelinks]{hyperref}

\newtheorem{theorem}{Theorem}
\newtheorem{corollary}[theorem]{Corollary}
\newtheorem{lemma}[theorem]{Lemma}
\newcommand{\R}{\mathbb R}
\newcommand{\E}{\mathbb E}
\newcommand{\tr}{\operatorname{tr}}
\newcommand{\ip}[2]{\left\langle #1,#2\right\rangle}

\title{An exact quadratic-chaos reduction for Heilman's Question 1.1\\
and the sign-paired spectral subcase}
\author{Automated mathematical discovery note (GPT5.6)}
\date{22 July 2026}

\begin{document}
\maketitle

\paragraph{Status.}
Likely valid partial result.  The general asymmetric-spectrum case remains
open in this note.

\section{Source question}

Let $\gamma_n$ be standard Gaussian measure on $\R^n$.  In the Introduction
of \cite{Heilman2015}, Question 1.1 asks, without a constraint on
$\gamma_n(A)$, whether
\[
 F(A)\leq \frac1\pi
\]
for every measurable symmetric set $A=-A$.  After rotating coordinates so
that the centered second-moment matrix is diagonal, the functional in the
paper is
\[
 F(A)=\sum_{i=1}^n\left(\int_A(1-x_i^2)\,d\gamma_n(x)\right)^2.
\]

Define the coordinate-free matrix
\[
 C_A:=\int_A(I-xx^{\mathsf T})\,d\gamma_n(x).
\]
Orthogonal invariance and diagonalization give $F(A)=\|C_A\|_F^2$, where
$\|\cdot\|_F$ denotes the Frobenius norm.

\section{Idea of the proof}

Frobenius duality turns $\|C_A\|_F$ into an integral of a centered Gaussian
quadratic form.  For a fixed quadratic form, its positive and negative sets
are themselves symmetric, and choosing one of them captures exactly half of
the form's first absolute moment.  This gives an exact equivalence, not only
a one-sided reduction.

For a spectrum made of opposite pairs, the characteristic function of the
quadratic form is real and positive.  The elementary inequality
$\log(1+s)\leq s$ compares it pointwise with the characteristic function
$e^{-t^2}$ of a Gaussian with variance two.  The standard Fourier formula
for $\E|Z|$ then proves the desired constant.

\section{Exact reduction}

\begin{lemma}[Best set for a centered even function]\label{lem:set}
Let $q\in L^1(\gamma_n)$ be even and satisfy $\int q\,d\gamma_n=0$.  Then
\[
 \sup_{A=-A}\left|\int_Aq\,d\gamma_n\right|
 =\frac12\int_{\R^n}|q|\,d\gamma_n.
\]
The supremum is attained, up to null sets, by $\{q>0\}$ and by $\{q<0\}$.
\end{lemma}

\begin{proof}
Write $q=q_+-q_-$.  The mean-zero assumption gives
$\int q_+\,d\gamma_n=\int q_-\,d\gamma_n=\frac12\int|q|\,d\gamma_n$.
For every measurable $A$,
\[
 -\int q_-\,d\gamma_n\leq \int_Aq\,d\gamma_n
 \leq \int q_+\,d\gamma_n.
\]
Since $q$ is even, both sign sets are symmetric and attain the endpoints.
\end{proof}

\begin{theorem}[Exact quadratic-chaos formulation]\label{thm:reduction}
If $X\sim N(0,I_n)$, then
\[
 \sup_{A=-A}\|C_A\|_F
 =\frac12\sup_{\substack{U=U^{\mathsf T}\\\|U\|_F=1}}
 \E\left|X^{\mathsf T}UX-\tr U\right|.
\]
Consequently, Heilman's Question 1.1 is equivalent to
\begin{equation}\label{eq:chaos}
 \E\left|X^{\mathsf T}UX-\tr U\right|\leq\frac{2}{\sqrt\pi}
 \qquad(U=U^{\mathsf T},\ \|U\|_F=1).
\end{equation}
For a fixed $U$, a maximizing set in the inner optimization is a quadratic
threshold set.
\end{theorem}

\begin{proof}
For symmetric matrices use the Frobenius inner product.  Put
$q_U(x)=x^{\mathsf T}Ux-\tr U$.  Then $q_U$ is even, Gaussian-integrable,
and centered, while
\[
 \ip{C_A}{U}_F
 =\int_A\tr\bigl(U(I-xx^{\mathsf T})\bigr)\,d\gamma_n(x)
 =-\int_Aq_U(x)\,d\gamma_n(x).
\]
Therefore, by finite-dimensional norm duality and Lemma~\ref{lem:set},
\begin{align*}
 \sup_{A=-A}\|C_A\|_F
 &=\sup_{A=-A}\sup_{\|U\|_F=1}|\ip{C_A}{U}_F|\\
 &=\sup_{\|U\|_F=1}\sup_{A=-A}\left|\int_Aq_U\,d\gamma_n\right|\\
 &=\frac12\sup_{\|U\|_F=1}\E|q_U(X)|.
\end{align*}
The sign sets of $q_U$ are the asserted quadratic thresholds.  Finally,
$F(A)=\|C_A\|_F^2$, so the bound $F(A)\leq1/\pi$ for every $A$ is
equivalent to \eqref{eq:chaos}.
\end{proof}

\section{The sign-paired spectral subcase}

\begin{theorem}[Opposite-pair spectra]\label{thm:pairs}
Let $U=U^{\mathsf T}$, $\|U\|_F=1$, and suppose that the eigenvalue multiset
of $U$ is
\[
 (a_1,-a_1,\ldots,a_m,-a_m,0,\ldots,0).
\]
Then
\[
 \E\left|X^{\mathsf T}UX-\tr U\right|<\frac{2}{\sqrt\pi}.
\]
In particular, the non-strict inequality in \eqref{eq:chaos} holds for this
entire spectral class.
\end{theorem}

\begin{proof}
By rotational invariance, with independent standard Gaussians $G_j,H_j$,
\[
 Q:=X^{\mathsf T}UX-\tr U
 \stackrel{d}{=}\sum_{j=1}^m a_j(G_j^2-H_j^2).
\]
The Frobenius normalization says $2\sum_ja_j^2=1$.  Independence gives the
real positive characteristic function
\[
 \phi_Q(t)=\prod_{j=1}^m(1+4a_j^2t^2)^{-1/2}.
\]
Since $\log(1+s)\leq s$ for $s\geq0$,
\[
 \log\phi_Q(t)
 =-\frac12\sum_j\log(1+4a_j^2t^2)
 \geq-2t^2\sum_ja_j^2=-t^2.
\]
Thus $\phi_Q(t)\geq e^{-t^2}$.  For a centered integrable random variable
with real characteristic function, the first absolute-moment identity is
\[
 \E|Q|=\frac2\pi\int_0^\infty\frac{1-\phi_Q(t)}{t^2}\,dt.
\]
It follows that
\[
 \E|Q|\leq\frac2\pi\int_0^\infty
 \frac{1-e^{-t^2}}{t^2}\,dt
 =\frac2{\sqrt\pi},
\]
where integration by parts gives the final integral as $\sqrt\pi$.
If $U\neq0$, then at least one $a_j$ is nonzero and
$\log(1+s)<s$ for $s>0$; hence the characteristic-function comparison is
strict for every $t\neq0$, which makes the final inequality strict.
\end{proof}

\begin{corollary}[A class of symmetric sets]\label{cor:sets}
Suppose the nonzero eigenvalues of $C_A$ occur in opposite pairs.  Then
$F(A)\leq1/\pi$, strictly if $C_A\neq0$.  In dimension two, this applies
whenever $\tr(C_A)=0$.
\end{corollary}

\begin{proof}
If $C_A=0$ there is nothing to prove.  Otherwise set
$U=C_A/\|C_A\|_F$.  Its spectrum is sign-paired, and
\[
 \|C_A\|_F=|\ip{C_A}{U}_F|
 \leq\frac12\E|q_U(X)|<\frac1{\sqrt\pi}
\]
by Lemma~\ref{lem:set} and Theorem~\ref{thm:pairs}.  Square the inequality.
For a two-by-two symmetric matrix, trace zero is equivalent to its
eigenvalues being $a,-a$.
\end{proof}

\section{Verification, limitations, and literature status}

The proof above is finite-dimensional and analytic; no numerical statement
is used.  A Monte Carlo subgradient search over dimensions up to twenty found
no violation of \eqref{eq:chaos}.  It approached $2/\sqrt\pi$ through diffuse
positive spectra, consistent with the central limit theorem, but this is only
supporting evidence.

For a general spectrum,
\[
 \phi_Q(t)=\prod_i e^{-it\lambda_i}(1-2it\lambda_i)^{-1/2}
\]
has a phase.  Its real part need not dominate $e^{-t^2}$ pointwise, so the
argument of Theorem~\ref{thm:pairs} does not extend directly.  This is the
precise remaining obstruction.

Bounded searches used the arXiv id, the exact wording of Question 1.1, and
phrases including ``second Wiener chaos $L^1$'', ``mean absolute deviation
Gaussian quadratic form'', and ``weighted centered chi-square absolute
moment''.  No paper explicitly resolving Question 1.1 was found.  However,
Theorem 3.1 of \cite{Klebanov2019} proves the same Gaussian first-moment upper
bound for every symmetric infinitely divisible distribution.  Since the
sign-paired $Q$ above is symmetric and infinitely divisible, that theorem
independently subsumes the moment inequality in
Theorem~\ref{thm:pairs}.  It does not state the exact set/quadratic-chaos
reduction or identify Heilman's matrix-functional subcase.  Novelty is
therefore limited and should be assessed by a specialist.

\begin{thebibliography}{9}
\bibitem{Heilman2015}
S. Heilman,
\emph{Low Correlation Noise Stability of Symmetric Sets},
arXiv:1511.00382 (2015).
\par\medskip

\bibitem{Klebanov2019}
L. B. Klebanov, A. V. Kakosyan, and I. V. Volchenkova,
\emph{Inequalities for $m$-Divisible Distributions and Testing of Infinite
Divisibility}, arXiv:1904.07604 (2019), Theorem 3.1.
\end{thebibliography}

\end{document}
